%% USPSC-Agradecimentos.tex
\begin{agradecimentos}

    Diante do complexo desafio envolvido no desenvolvimento desta dissertação de mestrado, permeado por inúmeros obstáculos pessoais e profissionais, encontro-me, ao término desta jornada, perante uma oportunidade singular para expressar minha profunda gratidão a todos que, de alguma forma, contribuíram para a realização deste trabalho. Que este seja um espaço de reconhecimento e celebração, onde se registre a importância de cada pessoa que esteve presente nessa trajetória.

    Dou início agradecendo à Deus. A Ele, toda a honra e toda a glória, por ter me concedido forças e sabedoria para superar os desafios que surgiram ao longo do caminho. Sou grato pela saúde, coragem e determinação que me permitiram seguir em frente, mesmo diante das adversidades. Agradeço pela oportunidade de aprendizado e crescimento contínuo, que me tornou uma pessoa mais forte e resiliente. Que Sua luz continue a me guiar e proteger, juntamente com a intercessão de Nossa Senhora, nossa Santa Mãe.

    Agradeço à minha grande companheira, incentivadora e querida esposa, Carolina. Você é minha inspiração, meu porto seguro, minha melhor amiga e o grande amor da minha vida. Obrigado por estar do meu lado em todos os momentos, por me apoiar, me encorajar e me motivar a seguir em frente, mesmo quando tudo parece difícil. Sua sensatez e lealdade são fundamentais para o meu desenvolvimento. Sou imensamente grato por sua presença, por iluminar nossa família com sua força e sabedoria e, especialmente, por me conceder a maior alegria da minha vida: nossa filha, Catarina.

    A você, minha filha, dedico estas palavras na tentativa de expressar o quão importante você é para nós. Enquanto escrevo, você está prestes a chegar ao mundo, e já sinto um amor incondicional por você. Espero que um dia leia este documento, seja para buscar uma referência ou por qualquer outro motivo... talvez tenha irmãos(ãs) ao seu lado nesse momento. Assim, aproveito para lhes dizer: nossa família é o bem mais precioso que temos. Papai e mamãe estarão sempre aqui por vocês. Eu amo vocês!

    O meu agradecimento se estende aos meus pais, Alessandra e Daniel, professores da minha vida e inspiradores da minha conduta. Obrigado por me ensinarem o valor da educação, do trabalho e da honestidade. Por apoiarem minhas escolhas e por acreditarem em meu potencial. Sou grato por suas cobranças, que sempre visaram meu melhor, e por sua confiança em mim. A vocês, consagro todo o meu esforço e dedicação. Ao meu irmão, Nícolas, agradeço por ser meu grande amigo e companheiro de vida. Desde que o segurei nos braços pela primeira vez, soube que se tornaria um homem íntegro e de coração generoso. Vocês são minha ligação com o passado e o alicerce do meu futuro.

    Registro um agradecimento especial ao Professor Ivan Nunes, meu orientador, por sua paciência, dedicação e orientação ao longo deste projeto. Agradeço por todo o conhecimento compartilhado, pelo incentivo à pesquisa e pela inspiração a seguir na carreira acadêmica. Sua contribuição foi essencial para o desenvolvimento desta dissertação.

    Ao Professor André Dias, meu mentor desde os primeiros anos da graduação, expresso minha profunda gratidão. Sua orientação constante nos projetos de pesquisa e inovação tecnológica foi fundamental para minha trajetória. Além de professor, tive a honra de tê-lo como amigo e conselheiro. Obrigado por acreditar no meu potencial e por me incentivar a sempre seguir em frente.

    Devo agradecer à empresa Nova Smar, pela oportunidade de atuar em um ambiente inovador e desafiador. Em especial, ao Diretor Libanio Souza e ao Gerente do departamento de P\&D, Delcio Prizon, pelo apoio incondicional à minha formação acadêmica e pelo incentivo à inovação. A confiança depositada em mim e as oportunidades concedidas foram essenciais para meu crescimento profissional. Também estendo minha gratidão a todos os colegas de trabalho que compartilharam essa jornada comigo.

    À Universidade de São Paulo (USP) e ao Instituto Federal de Educação, Ciência e Tecnologia de São Paulo (IFSP), por disponibilizarem os recursos necessários para o desenvolvimento deste. Um agradecimento especial ao colega Marcio Buzoli, por sua valiosa assistência na coleta de dados e na execução dos experimentos.

    Honro o fechamento deste ciclo expressando minha gratidão a todas as pessoas que, direta ou indiretamente, contribuíram para a concretização desta dissertação. Amigos, familiares, professores e colegas do meio acadêmico, corporativo e religioso. A convivência com pessoas tão especiais me motiva em cada etapa da jornada. Todo o percurso foi um aprendizado valioso e um exercício de gratidão. Tenho certeza de que, nos desafios futuros, estarei ainda mais preparado.
	
\end{agradecimentos}